\section{Gaussian contact as a Peano remainder}\label{gc:section}

The initial support bounds admit a single analytic proof. The key is to
recognize the contact kernel as a Gaussian Hermite remainder, and then
apply the classical Peano-kernel representation. This replaces the
initial actual-index support calculations and supplies the radial
contractions used by the tail route. Remaining actual-index refinements
and coordinate-strip exclusions retain their own proofs.

Let $Z$ be standard normal. For $g\in C^3(\mathbb R)$ with bounded
third derivative, put
\[
 L_xg=g(x)-\E g(Z)-x\E g'(Z)-\frac{x^2-1}{2}\E g''(Z).
\]
These expectations exist because $g,g',g''$ grow at most cubically,
quadratically and linearly.

\begin{lemma}[A uniform Gaussian Peano kernel]\label{gc:kernel}
For $|x|\ge2$,
\begin{equation}\label{gc:kernel-bound}
 |L_xg|\le\|g'''\|_\infty
       \left\{\frac{|x|^3-3|x|}{6}+\frac34\right\}.
\end{equation}
\end{lemma}
\begin{proof}
Taylor's integral remainder anchored at zero, followed by Fubini, gives
\begin{equation}\label{gc:peano}
 \begin{aligned}
 L_xg&=\int_{\mathbb R}g'''(t)K_x(t)\,dt,\\
 K_x(t)&=\frac{(x-t)^2}{2}
       \{\mathbf1_{\{t<x\}}-\overline\Phi(t)\}+(t/2-x)\phi(t).
 \end{aligned}
\end{equation}
Indeed $L_x$ annihilates all quadratic polynomials. Applying it to
$(u-t)_+^2/2$, and using the first two truncated Gaussian moments,
gives exactly this kernel. The anchored remainders are bounded by
Gaussian absolute moments of orders at most three, justifying the
interchanges. The displayed kernel has integrable Gaussian tails.
Testing $g(u)=u^3/6$ and reflecting give
\[
 \int K_x(t)\,dt=\frac{x^3-3x}{6},\qquad
 K_{-x}(-t)=-K_x(t)\quad\text{almost everywhere}.
\]

Take $x\ge2$. We first show $K_x(t)\ge0$ for $t\le x/2$.
For $t=-y\le0$, this is the Mills inequality
\[
 m(y):=\frac{\overline\Phi(y)}{\phi(y)}
 \ge\frac{2x+y}{(x+y)^2}.
\]
Cauchy--Schwarz for the integrals of $1,u,u^2$ against
$e^{-yu-u^2/2}du$ gives $m^2+ym\ge1$. Hence
\[
 m(y)\ge\frac{\sqrt{y^2+4}-y}{2}
       \ge\frac{y+4}{(y+2)^2}
       \ge\frac{y+2x}{(y+x)^2}.
\]
The middle comparison follows by squaring positive quantities; the
last follows by differentiation in $x\ge2$.

For $0\le t\le1$, $K_x(t)$ increases in $x\ge2$, since
$\partial_xK_x=(x-t)\Phi(t)-\phi(t)>0$. The elementary bounds
$39/100<\phi(0)<2/5$, $e^{-u}\ge1-u$ and
$e^{-u}\le1-u+u^2/2$ give $K_2(t)\ge P(t)$, where
\[
 P(t)=\frac15-\frac t{50}-\frac{13t^2}{100}
       -\frac{7t^3}{200}+\frac{3t^4}{100}-\frac{3t^5}{400}.
\]
On $[0,1]$, $P'(t)\le-49t/200-1/50<0$, so
$P(t)\ge P(1)=3/80>0$.
For $1\le t\le x/2$ the same monotonicity in $x$ gives
\[
 K_x(t)\ge K_{2t}(t)
       =\frac t2\{t\Phi(t)-3\phi(t)\}>0.
\]
The expression in braces increases for $t\ge1$; at $1$ it is
positive because $\Phi(1)\ge33/40$ and $\phi(1)<49/200$.

Writing $a=x/2$, we can now bound the negative part by
\begin{align}\label{gc:negative-part}
 \int K_x^-&\le\int_a^{2x}(x-t/2)\phi(t)\,dt
       +\frac12\int_x^\infty(t-x)^2\overline\Phi(t)\,dt\notag\\
 &\le\frac32\phi(a)+\frac12\phi(2x)+\frac{\phi(x)}{x^4}
 <\frac{5943}{16000}<\frac38.
\end{align}
The first integral uses $\overline\Phi(a)\le\phi(a)/a$.
The second is $\E(Z-x)_+^3/6$, at most $\phi(x)/x^4$ by
$\phi(x+u)\le\phi(x)e^{-xu}$. The last rational bound uses
$\phi(1)<49/200$, $\phi(2)<11/200$ and $\phi(4)<1/1000$.
Each follows from a positive partial sum of the exponential series.
Thus $\|K_x\|_1=\int K_x+2\int K_x^-$ proves
\eqref{gc:kernel-bound}.

\end{proof}

The signed moment needed for contact also has a support-free bound.
For $a=\E|X|>0$, $\beta=\E[X|X|]$ and $\rho=\E|X|^3$,
\begin{equation}\label{ut:initial-signed-moment}
 a\rho\ge1+\beta^2,\qquad a^2+\beta^2\le1,\qquad
 |\beta|\le\sqrt{\frac{\rho\sqrt{\rho^2+8}-\rho^2-2}{2}}.
\end{equation}
For the first inequality, apply Bessel's inequality to $|X|^{3/2}$
and the orthogonal functions $|X|^{1/2}$ and
$\operatorname{sgn}(X)|X|^{1/2}$. Their squared norms are $a$ and
their inner products with $|X|^{3/2}$ are $1,\beta$.
Cauchy--Schwarz for $X$ and $|X|-a$ proves the second inequality;
eliminating $a$ proves the last.

\begin{proposition}[Initial support without index subdivisions]\label{gc:initial-support}
Let $(P,H)$ be a standardized positive closed global maximizing pair
at index $n\ge25$, with $C=C_n\ge\CE$ and $C_{n-1}\le C$.
Then
\[
 \rho\le11/10\ \Longrightarrow\ \operatorname{supp}P\subset(-12/5,12/5),
\]
\[
 \rho\le6/5\ \Longrightarrow\ \operatorname{supp}P\subset(-13/5,13/5).
\]
No higher moment or defect bound is required.
\end{proposition}
\begin{proof}
For $g(u)=n^{3/2}\Phi((H-u)/\sqrt{n-1})$, Gaussian convolution gives
\begin{align*}
 \E g(Z)&=n^{3/2}\Phi(H/\sqrt n),\\
 \E g'(Z)&=-n\phi(H/\sqrt n),\qquad
 \E g''(Z)=-H\phi(H/\sqrt n).
\end{align*}
Thus $L_xg=\mathcal H_n(H,x)$ exactly. Also
$\|g'''\|_\infty=(n/(n-1))^{3/2}\phi(0)<213/500$ for $n\ge25$.
For $s_n=\sqrt{n/(n-1)}$ and $e_n=n(s_n-1)-1/2$,
\begin{equation}\label{gc:scale}
 0<e_n<\frac{2}{5n}\le\frac2{125}.
\end{equation}
Indeed, for $z=1/n\le1/25$,
\[
 (1-z)(1+z/2+2z^2/5)^2-1
 =z^2\{1/20-13z/20-6z^2/25-4z^3/25\}>0.
\]
The bracket decreases and is positive at $1/25$; omitting the
quadratic correction proves the lower bound.

By \eqref{ut:initial-signed-moment}, $\rho\le11/10$ gives
$|\beta|<127/500$, and $\rho\le6/5$ gives $|\beta|<44/125$.
For either $(\rho_u,b,R)=(11/10,127/500,12/5)$ or
$(6/5,44/125,13/5)$, put
\begin{align*}
 Q(r)&=-r^3+3\rho_ur^2/2+3br+2\rho_u/125,\\
 V(r)&=\frac{213}{500}\{(r^3-3r)/6+3/4\}
                                      +\frac{819}{2000}Q(r).
\end{align*}
At $r=R$, $Q<0$, $Q'<-8$ and $Q''<-11$; since $Q'''=-6$,
the multiplier of $C$ remains negative on the whole ray.
The contact inequality \eqref{loc:support-test} therefore has
right side at most $V(r)$ for $|x|=r\ge R$, using
$C\ge\CE>819/2000$ in the negative term. Direct substitution gives
\[
 V(12/5)=-278919/1250000,\qquad V(13/5)=-85993/1250000.
\]
In both rows $V'(R)<-5/2$, $V''(R)<-7/2$ and $V'''<-2$.
Hence $V<0$ on both entire rays, contradicting contact at any
support point there, including atoms and accumulation points.
\end{proof}

\begin{corollary}[All-index radial support at zero predecessor loss]\label{gc:tail-support}
Let $(P,H)$ be a standardized positive closed global maximizing pair
with $n\ge47$, $C=C_n\ge\CE$ and $C_{n-1}\le C$.
For each row below, the already established moment bounds
$\rho\le\rho_u$, $|\beta|\le b$ imply
$\operatorname{supp}P\subset(-R,R)$:
\[
\begin{array}{c|c|c|c}
 \rho_u&b&R&\text{contact upper bound at }R\\\hline
 7361/4000&643/1000&96/25&-1/20\\
 139/100&43/100&3&-23/100\\
 12563/10000&1963/5000&27/10&-9/500\\
 603687/500000&4014/15625&5/2&-49/500
\end{array}
\]
\end{corollary}
\begin{proof}
The predecessor allowance is zero. The same Gaussian function has
$\|g'''\|_\infty<207/500$ and $e_n<2/235$ for $n\ge47$.
For each row put
\begin{align*}
 Q(r)&=-r^3+3\rho_ur^2/2+3br+2\rho_u/235,\\
 V(r)&=\frac{207}{500}\{(r^3-3r)/6+3/4\}
       +\frac{819}{2000}Q(r).
\end{align*}
Direct substitution gives the strict upper bounds in the table.
In all four rows $Q(R)<0$, $Q'(R)<-8$, $Q''(R)<-11$,
$V'(R)<-5/2$, $V''(R)<-7/2$ and $V'''<-2$.
The whole-ray argument in Proposition~\ref{gc:initial-support}
therefore applies unchanged.
\end{proof}

These moment bounds are used in predecessor order. The global envelope
supplies the first row. The original central preparation, performed
at radius $3.84$, supplies the last row for all nine surviving central
intervals. The complementary acquisition at $3.84$ supplies the second
row; after contraction to $3$, its left-input and signed-resource
calculations supply the third. Thus no contraction uses a moment
bound obtained from its own output radius. The later coordinate-strip
exclusions retain their original bounded-threshold proofs.
The unsigned $\beta$ caps in these resource rows follow from
$\beta=(1-\theta)(1-2p)+\theta_+-\theta_-$, hence
$|\beta|\le(1-\theta)|1-2p|+\theta$; no sign of $\beta$ is assumed.
